\section{Q2 : expliquent-ils le comportement humain ?}
\frame{\sectionpage}

\begin{frame}{La surprise d'un mot}
\[
    s(w_t) = -\log p(w_t \mid w_1,\dots,w_{t-1})
\]
\begin{itemize}
    \item Un mot improbable est plus « surprenant »
    \item Les humains lisent plus lentement les mots surprenants
    \item Un modèle de langage fournit directement cette probabilité
\end{itemize}
\src{Levy, Cognition 2008 ; Smith \& Levy, Cognition 2013}
\end{frame}

\begin{frame}{Pourquoi les humains oublient parfois}
\begin{itemize}
    \item La mémoire est limitée : on ne garde pas tout le contexte
    \item Théorie de la \textbf{surprise à contexte dégradé} (lossy-context surprisal)
    \item Explique pourquoi les langues rapprochent les mots liés syntaxiquement
\end{itemize}
% TODO : vérifier les références de Futrell (2015 PNAS dependency length ; 2020 Cognitive Science)
\src{Futrell, Gibson \& Levy, Cognitive Science 2020}
\end{frame}

\begin{frame}{La surprise explique le N400}
\begin{itemize}
    \item Le \textbf{N400} : une onde EEG ($\sim$400 ms) qui augmente quand un mot est inattendu
    \item La surprise de modèles de langage explique de nombreux effets du N400
    \item Reste des cas où l'humain et le modèle divergent
\end{itemize}
% IMAGE : exemple d'onde N400 (à ajouter)
% TODO : vérifier la référence (Michaelov et al., Neurobiology of Language 2024)
\src{Michaelov, Bardolph, Van Petten, Bergen \& Coulson, 2024}
\end{frame}

\begin{frame}{Un meilleur modèle est-il plus humain ?}
\begin{itemize}
    \item Hypothèse admise en anglais : plus la perplexité est basse, plus le modèle prédit les temps de lecture
    \item En japonais, cette règle ne tient pas : \textbf{une perplexité plus basse n'est pas toujours plus humaine}
    \item Il faut évaluer les modèles sur plusieurs langues
\end{itemize}
\begin{block}{Message}
    Plus gros n'est pas automatiquement plus proche de l'humain.
\end{block}
\src{Kuribayashi, Oseki, Ito, Yoshida, Asahara \& Inui, ACL 2021}
\end{frame}

\begin{frame}{Comment le sens se compose}
\begin{itemize}
    \item « maison de chien » $\neq$ « chien de maison »
    \item Des modèles simples de composition sémantique sont testés sur des données fMRI
    \item Une composition apparaît dans le gyrus frontal inférieur gauche, avec un modèle explicite de la combinaison
\end{itemize}
\src{Ciapparelli, Marelli, Graves \& Reverberi, Cerebral Cortex 2025}
\end{frame}
